% !PRINTED 294
\weberpage{294}
%<<CONT
and then, if $\begin{pmatrix}\alpha, & \beta\\ \gamma, & \delta\end{pmatrix}$ is an arbitrary substitution in $\mathfrak{L}_0$ and
$\beta$ is different from $0$:
\weq{12}{\begin{pmatrix}\alpha, & \beta\\ \gamma, & \delta\end{pmatrix}=\begin{pmatrix}1, & 0\\ \delta\beta^{-1}, & 1\end{pmatrix}\begin{pmatrix}\beta, & 0\\ 0, & \beta^{-1}\end{pmatrix}\begin{pmatrix}0, & 1\\ -1, & 0\end{pmatrix}\begin{pmatrix}1, & 0\\ \alpha\beta^{-1}, & 1\end{pmatrix},}
and if $\beta=0$:
\weq{13}{\begin{pmatrix}\alpha, & 0\\ \gamma, & \alpha^{-1}\end{pmatrix}=\begin{pmatrix}\alpha, & 0\\ 0, & \alpha^{-1}\end{pmatrix}\begin{pmatrix}1, & 0\\ \gamma\alpha, & 1\end{pmatrix}.}
\gesperrt{From this theorem it is to be concluded that a group contained
 in $\mathfrak{L}_0$ which contains the two substitutions $A, B$
must necessarily be identical with $\mathfrak{L}_0$.}

\begin{center}
\S\,80. \gesperrt{Normal divisors of the group $\mathfrak{L}_0$.}
\end{center}

In order to learn the reductions of the transformation equation
 that may perhaps be possible, it is above all necessary to investigate the divisors
 of the group $\mathfrak{L}_0$. We ask first after
the existence of a \gesperrt{normal divisor} $\mathfrak{R}$ of $\mathfrak{L}_0$ (Bd. II, \S\,3)\footnote{After \gesperrt{Galois}: \glqq Proper divisors\grqq{}.}.

If
\weq{1}{S=\begin{pmatrix}\alpha, & \beta\\ \gamma, & \delta\end{pmatrix}}
is any substitution in $\mathfrak{R}$ different from the identical substitution,
 then, by the nature of the normal divisor, every
substitution
\weq{2}{U=TST^{-1}}
is likewise contained in $\mathfrak{R}$, if
\weq{3}{T=\begin{pmatrix}\alpha', & \beta'\\ \gamma', & \delta'\end{pmatrix}}
is an \gesperrt{arbitrary} substitution in $\mathfrak{L}_0$.
We set ourselves the task of determining $T$ and $\xi$ in such a way that
\weq{4}{U=\begin{pmatrix}0, & 1\\ -1, & \xi\end{pmatrix}}
results. This condition can also be written thus:
\[
\begin{pmatrix}\alpha', & \beta'\\ \gamma', & \delta'\end{pmatrix}\begin{pmatrix}\alpha, & \beta\\ \gamma, & \delta\end{pmatrix}=\begin{pmatrix}0, & 1\\ -1, & \xi\end{pmatrix}\begin{pmatrix}\alpha', & \beta'\\ \gamma', & \delta'\end{pmatrix},
\]
%CONT>>
